\documentclass{article}

\usepackage{amsmath, amsthm, amssymb, amsfonts}
\usepackage{thmtools}
\usepackage{graphicx}
\usepackage{setspace}
\usepackage{geometry}
\usepackage{float}
\usepackage[hidelinks]{hyperref}
\usepackage[utf8]{inputenc}
\usepackage[spanish]{babel}
\usepackage{framed}
\usepackage[dvipsnames]{xcolor}
\usepackage{tcolorbox}
\usepackage{tikz}
\usepackage{caption}
\usepackage{longtable}
\usepackage{pdflscape}
\usepackage{svg}
\usepackage{subcaption}
\usepackage{caption}
\usepackage{multirow}
\usepackage{array}
\usepackage{listings}
\usepackage{cancel}
\usepackage{fancyhdr}


\colorlet{LightGray}{White!90!Periwinkle}
\colorlet{LightOrange}{Orange!15}
\colorlet{LightGreen}{Green!15}



\newcommand{\HRule}[1]{\rule{\linewidth}{#1}}

\declaretheoremstyle[name=Theorem,]{thmsty}
\declaretheorem[style=thmsty,numberwithin=section]{theorem}
\tcolorboxenvironment{theorem}{colback=LightGray}

\declaretheoremstyle[name=Proposition,]{prosty}
\declaretheorem[style=prosty,numberlike=theorem]{proposition}
\tcolorboxenvironment{proposition}{colback=LightOrange}

\declaretheoremstyle[name=Principle,]{prcpsty}
\declaretheorem[style=prcpsty,numberlike=theorem]{principle}
\tcolorboxenvironment{principle}{colback=LightGreen}

\newcolumntype{L}[1]{>{\raggedleft\let\newline\\\arraybackslash\hspace{0pt}}m{#1}}
\newcolumntype{C}[1]{>{\centering\let\newline\\\arraybackslash\hspace{0pt}}m{#1}}
\newcolumntype{R}[1]{>{\raggedright\let\newline\\\arraybackslash\hspace{0pt}}m{#1}}

\setstretch{1.2}
\geometry{
    textheight=9in,
    textwidth=5.5in,
    top=1in,
    headheight=12pt,
    headsep=25pt,
    footskip=30pt
}

\lstdefinestyle{pythonstyle}{
    language=Python,
    basicstyle=\ttfamily,
    backgroundcolor=\color{gray!10},
    keywordstyle=\color{blue},
    commentstyle=\color{green!40!black},
    stringstyle=\color{red},
    showstringspaces=false,
    numbers=left,
    numberstyle=\tiny\color{gray},
    breaklines=true,
    breakatwhitespace=true,
    frame=tb,
    rulecolor=\color{black!70},
    framerule=0.5pt,
    tabsize=4,
    captionpos=b
}

% --------------------------------------------------------------------
% TITLE PAGE
% --------------------------------------------------------------------
\title{
    \vspace{2cm}
    \HRule{1pt} \\[0.5cm]
    \textbf{\Huge Plantilla de Programación Competitiva - Python}\\[0.3cm]
    \HRule{1pt} \\[1.5cm]
    \LARGE Universidad de Bogotá Jorge Tadeo Lozano \\[0.3cm]
    \Large Facultad de Ciencias Naturales e Ingeniería\\[0.3cm]
    \Large Semillero de Programación Competitiva\\[1.5cm]
    \large \textbf{Maratón Nacional ACIS REDIS 2025}\\[2cm]
}

\author{
    \Large Elaborado por: Nicolás Leitón
}

\date{\vspace{1cm} \Large Bogotá D.C., \today}


% --------------------------------------------------------------------
\begin{document}

\maketitle
\newpage

\tableofcontents

\thispagestyle{empty}
\newpage
\setcounter{page}{1}




\section{Leer Input}

\subsection{n inicial y luego casos}

\begin{lstlisting}[style=pythonstyle]
import sys
inputs = iter(sys.stdin.readlines())
TC = int (next(inputs))
for _ in range(TC):
  print(sum(map(int, next(inputs).split())))

\end{lstlisting}

\begin{lstlisting}
3
1 2
5 7
6 3
--- out
3
12
9
\end{lstlisting}

\subsection{Terminar cuando es 0 0}

\begin{lstlisting}[style=pythonstyle]
import sys
for line in sys.stdin.readlines():
  if line=='0 0\n': break
  print(sum(map(int, line.split())))
\end{lstlisting}

\begin{lstlisting}
1 2
5 7
6 3
0 0
1 1 <- no se procesa
--- out
3
12
9
\end{lstlisting}


\subsection{Enumerar casos hasta EOF}

\begin{lstlisting}[style=pythonstyle]

import  sys
for c, line in enumerate(sys.stdin.readlines(), 1):
    print("Case %s: %s\n" % (c, sum(map(int, line.split()))))

\end{lstlisting}


\begin{lstlisting}
1 2
5 7
6 3
--- out
Case 1: 3

Case 2: 12

Case 3: 9

\end{lstlisting}

\subsection{Input matriz nxm}

\begin{lstlisting}[style=pythonstyle]

import sys
lines = sys.stdin.readlines()
i= 0
matrices = []

while i<len(lines):
  # Leer dimensiones
  n, m = map(int, lines[i].split())
  i += 1

  #Leer la matriz de tamano n
  matrix = [list(map(int, lines[i+j].split())) for j in range(n)]
  matrices.append(matrix)

  i +=n

\end{lstlisting}




\section{Output Format}

\begin{lstlisting}[style=pythonstyle]
#Decimals output
text = "The price is {:.2f} dollars"
print(txt.format(45))
# The price is 45.00 dollar

#Use ',' as thousand separator:
txt = "the universe is {:,} years old "
print(txt.format(13800000000))
# The universe is 13,800,000,000 years old

# Convert to percentage format:
txt ="you scored {:.1%}"
print(txt.format(0.255)''
# you scored 25.5%

txt ="you scored {:.0%}"
print(txt.format(0.255)''
# you scored 26%


#covert number to binary

txt = "The binary version of {0} is {0:b}"
print(txt.format(5))
#The binary version of 5 is 101



\end{lstlisting}


\section{Macros}

\begin{lstlisting}[style=pythonstyle]
def crear_matriz(n, m, valor =0):
  return [[valor for _ in range(m)] for _ in range(n)]

def imprimir_matriz(mat):
  for fila in mat:
     print(*fila)

n, m = 3, 4
matriz= crear_matriz(n, m)
imprimir_matriz(matriz)

"""
0 0 0 0
0 0 0 0
0 0 0 0
"""

\end{lstlisting}


\section{Estructuras de Datos Built-In}

Esta sección contiene las estructuras de datos más útiles de Python para programación competitiva.

\subsection{Collections - Deque}

\begin{lstlisting}[style=pythonstyle]
from collections import deque

dq = deque([1, 2, 3])
dq.append(4)       # [1, 2, 3, 4]
dq.appendleft(0)   # [0, 1, 2, 3, 4]
dq.popleft()       # -> 0
\end{lstlisting}

\subsection{Heapq - Priority Queue}

\begin{lstlisting}[style=pythonstyle]
import heapq

h = []
heapq.heappush(h, 3)
heapq.heappush(h, 1)
heapq.heappush(h, 5)

print(heapq.heappop(h))  # 1 (menor elemento)

# Max heap:
h = []
heapq.heappush(h, -3)
heapq.heappush(h, -1)
heapq.heappush(h, -5)
print(-heapq.heappop(h)) # 5 (mayor)
\end{lstlisting}

\subsection{DefaultDict - Grafos}

\begin{lstlisting}[style=pythonstyle]
from collections import defaultdict
#Representar grafos con listas de adyacencia.
g = defaultdict(list)
n, m = 5, 6 #Nodos, aristas
edges = [(1,2), (1,3), (2,4), (3,4), (4,5), (5,1)]

for u, v in edges:
    g[u].append(v)
    g[v].append(u)  # si es no dirigido
\end{lstlisting}

\subsection{Counter - Frecuencias}

\begin{lstlisting}[style=pythonstyle]
from collections import Counter
#Cuenta frecuencias mas rapido
cnt = Counter([1,2,2,3,3,3])
print(cnt[2])       # 2
print(cnt.most_common(1))  # [(3, 3)]
\end{lstlisting}

\subsection{Bisect - Búsqueda Binaria}

\begin{lstlisting}[style=pythonstyle]
import bisect
#Insercion en lista ordenada
arr = [1, 3, 3, 5, 8]

print(bisect.bisect_left(arr, 3))   # -> 1
print(bisect.bisect_right(arr, 3))  # -> 3

#Cuantos elemento hay entre L y R
arr = [1, 2, 4, 4, 5, 7, 9]
L, R = 4, 7
left = bisect.bisect_left(arr, L)
right = bisect.bisect_right(arr, R)
print(right - left)  # 4 elementos en el rango [4,7]
\end{lstlisting}


\section{Algoritmos}

\subsection{Two Sum}

\begin{lstlisting}[style=pythonstyle]
#La funcion intenta determinar si en el arreglo a (de tamano n) 
#existen dos numeros cuya suma sea igual a k.
def isPairSum(a, n, k):
	global x, y
	i, j = 0, n - 1
	while i < j:
		if a[i] + a[j] == k:
			x, y = i, j
			return 1
		elif a[i] + a[j] > k:
			j = j - 1
		else:
			i = i + 1
	return 0
\end{lstlisting}

\subsection{Sliding Window - Suma Mínima}

\begin{lstlisting}[style=pythonstyle]
# Halla la suma minima de un subarreglo de k elementos en un array,
#  y devuelve los indices donde empieza y termina dicho subarreglo.
for _ in range(int(input())):
	n, k = map(int, input().split())
	arr = list(map(int, input().split()))
	cursum = res_ind = 0
	# cursum = sum(x for x in arr[:k])
	for i in range(k):
		cursum += arr[i]
	minsum = cursum
	for i in range(k, n):
		cursum = cursum + arr[i] - arr[i - k] #Deslizar ventana
		if cursum < minsum:
			minsum = cursum
			res_ind = i - k + 1
	print(res_ind, res_ind + k - 1)

'''
#----INPUT----

1
7 3
3 7 90 20 10 50 40

#----OUTPUT----

3 5
'''
\end{lstlisting}


\section{Grafos: Dijkstra}

\subsection{Dijkstra - Camino más corto}

Complejidad: \(O((N + M) \log N)\) con heapq.

\begin{lstlisting}[style=pythonstyle]
from heapq import heappush, heappop
import math

def dijkstra(graph, start):
    dist = {u: math.inf for u in graph}
    dist[start] = 0
    pq = [(0, start)]
    while pq:
        d, u = heappop(pq)
        if d > dist[u]: continue
        for v, w in graph[u]:
            if dist[v] > d + w:
                dist[v] = d + w
                heappush(pq, (dist[v], v))
    return dist

# Ejemplo de uso
graph = {
    1: [(2, 4), (3, 1)],
    2: [(3, 2), (4, 5)],
    3: [(4, 8)],
    4: []
}
#Distancias minimas desde start a cada nodo
distancias = dijkstra(graph, 1)
print(distancias) #{1: 0, 2: 3, 3: 1, 4: 8}
\end{lstlisting}

\textbf{Trucos:}
\begin{itemize}
  \item Usar heapq para min-heap eficiente.
  \item Para max-heap, negar los valores.
  \item Mantener diccionario de padres para reconstruir camino.
\end{itemize}


\section{Template Completo - Template.py}

Este es el template completo que incluye utilidades comunes para programación competitiva:

\begin{lstlisting}[style=pythonstyle]

# PRINT ARRAY WITHOUT [,,,]
#print(*arr)


#La funcion intenta determinar si en el arreglo a (de tamano n) 
#existen dos numeros cuya suma sea igual a k.
def isPairSum(a, n, k):
	global x, y
	i, j = 0, n - 1
	while i < j:
		if a[i] + a[j] == k:
			x, y = i, j
			return 1
		elif a[i] + a[j] > k:
			j = j - 1
		else:
			i = i + 1
	return 0


# Halla la suma minima de un subarreglo de k elementos en un array,
#  y devuelve los indices donde empieza y termina dicho subarreglo.
for _ in range(int(input())):
	n, k = map(int, input().split())
	arr = list(map(int, input().split()))
	cursum = res_ind = 0
	# cursum = sum(x for x in arr[:k])
	for i in range(k):
		cursum += arr[i]
	minsum = cursum
	for i in range(k, n):
		cursum = cursum + arr[i] - arr[i - k] #Deslizar ventana
		if cursum < minsum:
			minsum = cursum
			res_ind = i - k + 1
	print(res_ind, res_ind + k - 1)

'''
#----INPUT----

1
7 3
3 7 90 20 10 50 40

#----OUTPUT----

3 5
'''




'''
COLLECTIONS PYTHON
'''

from collections import deque

dq = deque([1, 2, 3])
dq.append(4)       # [1, 2, 3, 4]
dq.appendleft(0)   # [0, 1, 2, 3, 4]
dq.popleft()       # -> 0



import heapq

h = []
heapq.heappush(h, 3)
heapq.heappush(h, 1)
heapq.heappush(h, 5)

print(heapq.heappop(h))  # 1 (menor elemento)

# Max heap:
h = []
heapq.heappush(h, -3)
heapq.heappush(h, -1)
heapq.heappush(h, -5)
print(-heapq.heappop(h)) # 5 (mayor)




from collections import defaultdict
#Representar grafos con listas de adyacencia.
g = defaultdict(list)
n, m = 5, 6 #Nodos, aristas
edges = [(1,2), (1,3), (2,4), (3,4), (4,5), (5,1)]

for u, v in edges:
    g[u].append(v)
    g[v].append(u)  # si es no dirigido


from collections import Counter
#Cuenta frecuencias mas rapido
cnt = Counter([1,2,2,3,3,3])
print(cnt[2])       # 2
print(cnt.most_common(1))  # [(3, 3)]



import bisect
#Insercion en lista ordenada
arr = [1, 3, 3, 5, 8]

print(bisect.bisect_left(arr, 3))   # -> 1
print(bisect.bisect_right(arr, 3))  # -> 3

#Cuantos elemento hay entre L y R
arr = [1, 2, 4, 4, 5, 7, 9]
L, R = 4, 7
left = bisect.bisect_left(arr, L)
right = bisect.bisect_right(arr, R)
print(right - left)  # 4 elementos en el rango [4,7]



'''Djistra'''
from heapq import heappush, heappop
import math

def dijkstra(graph, start):
    dist = {u: math.inf for u in graph}
    dist[start] = 0
    pq = [(0, start)]
    while pq:
        d, u = heappop(pq)
        if d > dist[u]: continue
        for v, w in graph[u]:
            if dist[v] > d + w:
                dist[v] = d + w
                heappush(pq, (dist[v], v))
    return dist
graph = {
    1: [(2, 4), (3, 1)],
    2: [(3, 2), (4, 5)],
    3: [(4, 8)],
    4: []
}
#Distancias minimas desde start a cada nodo
distancias = dijkstra(graph, 1)
print(distancias) #{1: 0, 2: 3, 3: 1, 4: 8}

\end{lstlisting}


\section{Consejos Finales}

\begin{itemize}
  \item Usar \texttt{sys.stdin.readlines()} para lectura rápida.
  \item Preferir list comprehensions para mejor rendimiento.
  \item Usar \texttt{collections} para estructuras de datos eficientes.
  \item Familiarizarse con \texttt{heapq}, \texttt{bisect} y \texttt{itertools}.
  \item Para problemas de grafos, considerar usar diccionarios de adyacencia.
\end{itemize}

\end{document}
